\documentclass[10pt,a4paper]{article}

% ---- Packages ----
\usepackage[utf8]{inputenc}
\usepackage[T1]{fontenc}
\usepackage{lmodern}
\usepackage{microtype}
\usepackage[margin=1in]{geometry}
\usepackage{amsmath,amssymb}
\usepackage{booktabs}
\usepackage{graphicx}
\usepackage{url}
\usepackage{hyperref}
\usepackage{natbib}
\usepackage{caption}
\usepackage{subcaption}
\usepackage{xcolor}
\usepackage{enumitem}
\setlength{\emergencystretch}{1.5em} % third-pass stretch for paragraphs dense with code identifiers
% No algorithm2e / algorithm / algpseudocode assumed available in the
% target TeX distribution — we inline the procedure as a numbered list.

\hypersetup{colorlinks=true, linkcolor=blue!50!black, citecolor=blue!50!black,
            urlcolor=blue!50!black}

\graphicspath{{figures/}}

% ---- Result values ----
% Every result value is a macro generated from artifacts/v0.6.0 by
% scripts/paper/fill_v060_numbers.py (no number is typed by hand).
% GENERATED by scripts/paper/fill_v060_numbers.py from artifacts/v0.6.0 — do not edit by hand.
% Every result value in the manuscript is one of these macros (CTO #491: no hand-typed numbers).
\newcommand{\resArchiveShaCache}{703916525d4591ac}
\newcommand{\resArchiveShaOutputs}{419daa10f1492436}
\newcommand{\resCacheHits}{0}
\newcommand{\resCacheStart}{0}
\newcommand{\resCeilingArchitect}{256}
\newcommand{\resCeilingCurator}{256}
\newcommand{\resCeilingLine}{30}
\newcommand{\resCeilingLiterature}{1316}
\newcommand{\resCeilingTrainer}{256}
\newcommand{\resCeilingValidator}{1192}
\newcommand{\resGPUHours}{3.07}
\newcommand{\resGateSummary}{4 of 5 gates pass; H5 fails (0 unevaluated)}
\newcommand{\resGitDirty}{false}
\newcommand{\resGitSha}{ce5f237}
\newcommand{\resHFiveCIHi}{0.726}
\newcommand{\resHFiveCILo}{-0.244}
\newcommand{\resHFiveN}{20}
\newcommand{\resHFiveRho}{0.286}
\newcommand{\resHFiveThreshold}{0.4}
\newcommand{\resHFiveVerdict}{\textbf{FAIL}}
\newcommand{\resHFiveWAce}{-1.53}
\newcommand{\resHFiveWDc}{-0.29}
\newcommand{\resHFourAdaAceCIHi}{-0.182}
\newcommand{\resHFourAdaAceCILo}{-0.888}
\newcommand{\resHFourAdaAceN}{21}
\newcommand{\resHFourAdaAcePass}{no}
\newcommand{\resHFourAdaAceRho}{-0.628}
\newcommand{\resHFourAdaDcCIHi}{0.906}
\newcommand{\resHFourAdaDcCILo}{0.238}
\newcommand{\resHFourAdaDcN}{21}
\newcommand{\resHFourAdaDcPass}{yes}
\newcommand{\resHFourAdaDcRho}{0.663}
\newcommand{\resHFourAdaTdiCIHi}{0.904}
\newcommand{\resHFourAdaTdiCILo}{0.236}
\newcommand{\resHFourAdaTdiN}{21}
\newcommand{\resHFourAdaTdiPass}{yes}
\newcommand{\resHFourAdaTdiRho}{0.661}
\newcommand{\resHFourNPassing}{4}
\newcommand{\resHFourNTests}{6}
\newcommand{\resHFourNorAceCIHi}{0.141}
\newcommand{\resHFourNorAceCILo}{-0.794}
\newcommand{\resHFourNorAceN}{20}
\newcommand{\resHFourNorAcePass}{no}
\newcommand{\resHFourNorAceRho}{-0.376}
\newcommand{\resHFourNorDcCIHi}{0.891}
\newcommand{\resHFourNorDcCILo}{0.338}
\newcommand{\resHFourNorDcN}{20}
\newcommand{\resHFourNorDcPass}{yes}
\newcommand{\resHFourNorDcRho}{0.703}
\newcommand{\resHFourNorTdiCIHi}{0.889}
\newcommand{\resHFourNorTdiCILo}{0.337}
\newcommand{\resHFourNorTdiN}{20}
\newcommand{\resHFourNorTdiPass}{yes}
\newcommand{\resHFourNorTdiRho}{0.702}
\newcommand{\resHFourPooledAceRho}{-0.585}
\newcommand{\resHFourPooledDcRho}{0.769}
\newcommand{\resHFourPooledTdiRho}{0.767}
\newcommand{\resHFourRunsUndefined}{0}
\newcommand{\resHFourTasksDropped}{0}
\newcommand{\resHFourValue}{0.703}
\newcommand{\resHFourVerdict}{\textbf{PASS}}
\newcommand{\resHOneCIHi}{0.149}
\newcommand{\resHOneCILo}{0.111}
\newcommand{\resHOneFracOver}{0.14}
\newcommand{\resHOneIQRHi}{0.151}
\newcommand{\resHOneIQRLo}{0.089}
\newcommand{\resHOneMax}{7.821}
\newcommand{\resHOneMedian}{0.132}
\newcommand{\resHOneN}{21}
\newcommand{\resHOneThreshold}{0.20}
\newcommand{\resHOneVerdict}{\textbf{PASS}}
\newcommand{\resHThreeCeiling}{1.099}
\newcommand{\resHThreeEntropy}{0.859}
\newcommand{\resHThreeExecNeStated}{0}
\newcommand{\resHThreeNModels}{1}
\newcommand{\resHThreeNSteps}{369}
\newcommand{\resHThreePickLinear}{19}
\newcommand{\resHThreePickMlp}{167}
\newcommand{\resHThreePickScgpt}{183}
\newcommand{\resHThreeVerdict}{\textbf{PASS}}
\newcommand{\resHTwoCIHi}{0.395}
\newcommand{\resHTwoCILo}{0.092}
\newcommand{\resHTwoDoubletMedian}{0.517}
\newcommand{\resHTwoDoubletN}{5}
\newcommand{\resHTwoFracOver}{0.45}
\newcommand{\resHTwoIQRHi}{0.417}
\newcommand{\resHTwoIQRLo}{0.082}
\newcommand{\resHTwoMax}{0.729}
\newcommand{\resHTwoMedian}{0.184}
\newcommand{\resHTwoN}{20}
\newcommand{\resHTwoSingletonMedian}{0.149}
\newcommand{\resHTwoSingletonN}{15}
\newcommand{\resHTwoThreshold}{0.30}
\newcommand{\resHTwoVerdict}{\textbf{PASS}}
\newcommand{\resKillLine}{28}
\newcommand{\resManifestSha}{323af965b631ce83}
\newcommand{\resNCalls}{1846}
\newcommand{\resNDistinctConfigs}{7}
\newcommand{\resNDistinctFits}{19}
\newcommand{\resNFallback}{0}
\newcommand{\resNLifecycleRuns}{123}
\newcommand{\resNSteps}{1845}
\newcommand{\resNTrainerRecordsPerTask}{27}
\newcommand{\resNTrainerRuns}{1107}
\newcommand{\resPreregVersion}{v0.6.0-a4}
\newcommand{\resProjGPU}{2.90}
\newcommand{\resProjLLM}{2.93}
\newcommand{\resProjLatency}{9.2}
\newcommand{\resProjTotal}{7.19}
\newcommand{\resRefusals}{0}
\newcommand{\resReplay}{false}
\newcommand{\resRosterHaiku}{claude-haiku-4-5-20251001}
\newcommand{\resRosterSonnet}{claude-sonnet-5-5}
\newcommand{\resRunId}{20260929T035447Z-ce5f237}
\newcommand{\resServedMismatch}{0}
\newcommand{\resSpendGPU}{4.05}
\newcommand{\resSpendHaiku}{1.75}
\newcommand{\resSpendLLM}{3.52}
\newcommand{\resSpendPreflight}{0.0105}
\newcommand{\resSpendPrior}{1.3548}
\newcommand{\resSpendSonnet}{1.76}
\newcommand{\resSpendStopLine}{12}
\newcommand{\resSpendTotal}{8.93}
\newcommand{\resStopEndTurn}{1845}
\newcommand{\resStopMaxTokens}{1}
\newcommand{\resStopReasonNone}{none}
\newcommand{\resTallyFail}{1}
\newcommand{\resTallyOut}{5}
\newcommand{\resTallyPass}{4}
\newcommand{\resTallyUneval}{0}
\newcommand{\resTokensIn}{1,242,639}
\newcommand{\resTokensOut}{277,906}


% ---- Title ----
\title{Does Agent Confidence Entropy Predict Task Difficulty?\\
       A Pre-registered, Provenance-Complete Test of Agentic
       Hyperparameter Tuning for Perturb-seq Response Prediction}
\hypersetup{pdftitle={Does Agent Confidence Entropy Predict Task Difficulty?
  A Pre-registered, Provenance-Complete Test of Agentic Hyperparameter
  Tuning for Perturb-seq Response Prediction}}

\author{%
  Mangyin Mo\\
  Carnegie Mellon University\\
  \texttt{mangyinm@alumni.cmu.edu}\\
  ORCID: 0009-0009-5233-3142%
}

\date{}

% ---- Document ----
\begin{document}

\maketitle

\begin{abstract}
Multi-agent group-generation systems spend test-time compute uniformly
across tasks. We ask whether the agents' own confidence-and-critique trace
predicts how hard a task is, which would let the compute be allocated
adaptively. Taking as testbed a five-agent orchestrator for Perturb-seq
response prediction inspired by \textsc{CellForge}~\citep{chen2025cellforge}
and operated via \textsc{MassGen}~\citep{massgen2025github}, we define four
compute-free trace metrics --- Agent Confidence Entropy (ACE), Critique
Severity Dispersion (CSD), a convergence signal ($\Delta C$), and a Winner
Flip Rate (WFR) --- and their composite, the Task Difficulty Index (TDI).
We pre-register five hypotheses with gates: attainable error on
\textsc{Adamson 2016}~\citep{adamson2016perturbseq} (CRISPR interference) and
\textsc{Norman 2019}~\citep{norman2019perturbseq} (CRISPR activation), the
Architect agent's backbone-choice entropy, the rank correlation of each TDI
component with held-out error, and cross-dataset transfer of a calibrated
TDI. The design is $41$ held-out K562 tasks ($21$ Adamson, $15$ Norman
singletons, $5$ Norman doublets) drawn deterministically, a trainer sweep of
\resNTrainerRecordsPerTask{} records per task and three fixed-length agentic
lifecycle runs per task, all in one process. The LLM condition is a pinned
roster of Anthropic Claude models (Haiku 4.5 for the four proposer roles,
Sonnet 5.5 for the Validator) whose served model is recorded and asserted
per call; a run with any rule-based fallback step is invalid. Every run records its git revision, resolved
configuration, library versions, dataset SHA-256 digests, task plan and
label contract. Gate outcomes: \resGateSummary{} --- the trainer sweep attains the pre-registered
error on both screens (H1, H2), the Architect exercises its backbone choice
(H3), and the confidence-drop component and the composite rank tasks by
held-out error within both screens (H4); the Adamson-calibrated composite
does not transfer to Norman at the pre-registered level (H5, $\rho=\resHFiveRho$
against $>\resHFiveThreshold$). The confidence-entropy point estimates run against the pre-registered
direction in both screens (exploratory; no claim is made).
The values reported in the v0.5.0 release of this work are superseded in
full; the defects behind each are listed in the appendix.
\end{abstract}

% =====================================================================
\section{Introduction}\label{sec:intro}

Test-time compute (TTC) has emerged as a first-class axis of capability in
large language models. \citet{snell2024scaling} showed that an
\emph{optimally allocated} per-prompt TTC budget lets a small model match a
$14\times$ larger model under FLOPs-matched comparison, with the optimal
allocation depending on the difficulty of the prompt. The recent taxonomy of
\citet{zhang2025ttssurvey} organises the exploding TTC literature along four
orthogonal axes --- \emph{what} to scale, \emph{how} to scale, \emph{where} to
scale, and \emph{how well} the method scales --- but the prominent
applications to date address \emph{single-chain} reasoning: parallel samples
scored by a process reward model (PRM)~\citep{lightman2024prm}, sequential
self-revision, tree search~\citep{yao2023tot}, or internal chain-of-thought
trained via reinforcement learning~\citep{deepseek2025r1,muennighoff2025s1}.

The multi-agent case is structurally different but no less TTC-laden. A
group-generation system such as \textsc{CellForge}~\citep{chen2025cellforge}
dispatches several specialised agents --- typically a data curator, a
literature/prior agent, a model architect, a trainer, and a biological
validator --- and coordinates them through propose-critique-vote rounds. Every
such system today spends compute \emph{uniformly}: the same number of agents
for the same number of rounds regardless of whether the task is a
well-characterised single-gene perturbation or a novel combinatorial one.
\emph{This paper asks whether the group's own internal dynamics carry enough
signal to allocate its compute adaptively}, and tests that question under a
pre-registered design so that either answer is informative.

\paragraph{Contributions.} We make four contributions, all within the
empirical frame of Perturb-seq response prediction~\citep{norman2019perturbseq,%
adamson2016perturbseq} but generalising to any propose-critique-vote system:

\begin{enumerate}[leftmargin=*]
  \item We formalise four easy-to-compute observables of a multi-agent run ---
        ACE, CSD, $\Delta C$/$\Delta$ACE, and WFR --- and combine them into a
        scalar Task Difficulty Index (TDI).
  \item We introduce a \emph{preflight probe}: a single shallow round whose
        four-dimensional signature is intended to serve as a sufficient
        statistic for a Bayesian recommender over orchestration
        hyperparameters $\varphi = (N, R, \textrm{backbone})$.
  \item We pre-register five hypotheses and their gates
        (\S\ref{sec:results}; \texttt{paper/PREREGISTRATION.md}) and test them
        on $41$ held-out K562 tasks from \textsc{Adamson 2016} and
        \textsc{Norman 2019}, with a trainer sweep and an agentic lifecycle
        run over the identical task set in one process.
  \item We release a provenance-complete run record --- git revision,
        resolved configuration, resolved library versions, dataset digests,
        task plan, label contract, and the serving model of every agent step
        --- together with the calibration harness and the
        \textsc{MassGen}-installable evaluation skill
        (\S\ref{sec:reproducibility}).
\end{enumerate}

\paragraph{Scope.} Both datasets are K562 Perturb-seq screens from the
scPerturb~\citep{peidli2024scperturb} repackaging (Zenodo record
$13350497$). Genome-scale validation on
\textsc{Replogle 2022}~\citep{replogle2022perturbseq} is the natural next step
and is enumerated in \S\ref{sec:limits}; the paper does not extrapolate
beyond K562.

% =====================================================================
\section{Background and Related Work}\label{sec:related}

\paragraph{Test-time compute scaling.} The modern TTC literature is usefully
organised by whether the extra compute is spent \emph{in parallel} (best-of-$N$
with a verifier~\citep{cobbe2021verifiers,lightman2024prm}, self-consistency
voting~\citep{wang2023selfconsistency}), \emph{sequentially} (iterative
revision, self-refine), \emph{hybrid} (tree search, MCTS,
\citep{yao2023tot}), or \emph{internal} to the model
(\citep{deepseek2025r1,muennighoff2025s1}); see \citet{zhang2025ttssurvey} for
a taxonomy and \citet{wu2024smalltllm,liu2025rethinkingtts} for scaling-law
analyses. \citet{snell2024scaling} is the primary reference for the
difficulty-conditional compute-optimal idea that our work extends to the
group setting.

\paragraph{Multi-agent TTC.} \textsc{AutoGen}~\citep{wu2023autogen} and
\textsc{MassGen}~\citep{massgen2025github} generalise propose-critique-vote to
heterogeneous agent pools; \textsc{Archon}~\citep{saad2024archon} and
\textsc{MAV}~\citep{lifshitz2025mav} treat inference-time architecture
\emph{itself} as the search space. \textsc{CellForge}~\citep{chen2025cellforge}
applies this pattern specifically to single-cell experimental design. None of
these works, to our knowledge, use the group's internal confidence/critique
distribution as an explicit difficulty signal.

\paragraph{Perturb-seq and single-cell foundation models.} Perturb-seq
couples CRISPR perturbations with single-cell RNA-sequencing readouts,
offering a clean testbed for agentic scientific reasoning because ground-truth
hardness can be proxied from literature density, replicate agreement, and
donor-split transfer~\citep{norman2019perturbseq,adamson2016perturbseq,
replogle2022perturbseq}. Pretrained single-cell foundation models such as
\textsc{scGPT}~\citep{cui2024scgpt}, \textsc{Geneformer}%
~\citep{theodoris2023geneformer}, \textsc{scFoundation}~\citep{hao2024scfoundation},
and \textsc{scPRINT}~\citep{kalfon2025scprint} are related work; none of them
is among the backbones evaluated here (\S\ref{subsec:backbones}).

% =====================================================================
\section{Problem Setup}\label{sec:setup}

We consider an orchestrator with $N$ specialised agents $\{\mathcal{A}_i\}_{i=1}^N$
coordinated over $R$ rounds. At each round $r$, every agent produces a
structured proposal with scalar confidence $c_i(r) \in [0,1]$ and every other
agent emits a critique with severity $S_{ij}(r) \in [0,1]$ (row $i$ critic,
column $j$ target; self-critique is undefined). A winner index $w(r)$ is
assigned by the orchestrator (typically $\mathrm{argmax}\, c - \overline{S}$,
where $\overline{S}$ is the mean severity received). We let
$\mathbf{C}(r)\!\in\!\mathbb{R}^N$ and $\mathbf{S}(r)\!\in\!\mathbb{R}^{N\times(N-1)}$
denote the per-round tuples.

A task $t$ has a latent intrinsic difficulty $d_t\!\in\![0,1]$. The team's
downstream utility on the task is proxied by the prediction error
$u(t, \varphi)$ on the held-out perturbation (MSD on its top-$20$
differentially expressed genes, \S\ref{subsec:eval}), where
$\varphi = (N, R, b)$ with $b$ a backbone choice.
\emph{Agentic hyperparameter tuning} is the problem of picking
$\varphi^\star\!\in\!\Phi$ minimising compute subject to a target error ---
the multi-agent analogue of \citet{snell2024scaling}'s compute-optimal
single-chain allocation.

% =====================================================================
\section{Metrics}\label{sec:metrics}

\subsection{Agent Confidence Entropy (ACE)}

Let $p(r) = \mathrm{softmax}(\mathbf{C}(r)/\tau)$ be the softmax of the
confidence vector at round $r$ with temperature $\tau$ (we use $\tau=1$). ACE is the Shannon entropy of this distribution, normalised to
$[0,1]$ by its maximum:
\begin{equation}
  \mathrm{ACE}(r) = -\sum_{i=1}^{N} p_i(r) \log p_i(r),\quad
  \mathrm{ACE}_{\mathrm{norm}}(r) = \mathrm{ACE}(r)\,/\,\log N.
\end{equation}
$\mathrm{ACE}_{\mathrm{norm}}\!\to\!0$ when one agent dominates (a crisp hierarchy
among agents) and $\to\!1$ when the distribution is flat. Crucially, a flat
distribution has two very different causes --- \emph{confident consensus}
(easy task) and \emph{uniform uncertainty} (hard task) --- so ACE must be read
alongside the mean confidence.

\subsection{Critique Severity Dispersion (CSD)}

Let $\mathbf{S}(r)$ collect the off-diagonal entries of the critique matrix.
We define CSD as the population variance of its entries:
\begin{equation}
  \mathrm{CSD}(r) = \mathrm{Var}\left(\mathbf{S}(r)\right).
\end{equation}
Large CSD concentrates disagreement: a small number of severe critiques weighs
more than uniformly moderate ones. We also track the row- and column-variances
for diagnostic use (which critic is inconsistent? which proposal is contested?).

\subsection{Convergence signals}

For a run with $R$ rounds we track three convergence signals:
\begin{align}
  \Delta\mathrm{ACE}     &\triangleq \mathrm{ACE}_{\mathrm{norm}}(R) - \mathrm{ACE}_{\mathrm{norm}}(0), \\
  \Delta C               &\triangleq \overline{\mathbf{C}}(R) - \overline{\mathbf{C}}(0), \\
  \mathrm{WFR}           &\triangleq \frac{1}{R-1} \sum_{r=2}^{R} \mathbb{1}\!\left[w(r) \neq w(r-1)\right].
\end{align}
A well-posed run has $\Delta\mathrm{ACE}\!<\!0$, $\Delta C\!>\!0$, and
$\mathrm{WFR}\!\approx\!0$. We use $1-\min(\max(\Delta C,0),1)$ as a
bounded ``lack of convergence'' feature.

\subsection{Task Difficulty Index}

We combine the four signals linearly:
\begin{equation}
  \mathrm{TDI} = \alpha\,\mathrm{ACE}_{\mathrm{norm}}(R)
               + \beta\,\mathrm{CSD}(R)
               + \gamma\,\bigl(1 - \min(\max(\Delta C, 0), 1)\bigr)
               + \delta\,\mathrm{WFR},
  \label{eq:tdi}
\end{equation}
with non-negative weights that sum to unity so $\mathrm{TDI}\!\in\![0,1]$. We
compare two weight regimes: the hand-set \emph{default}
$(\alpha,\beta,\gamma,\delta)=(0.35,0.25,0.25,0.15)$ and the \emph{calibrated}
weights fitted by ridge regression of~(\ref{eq:tdi}) against a ground-truth
difficulty label.

% =====================================================================
\section{Preflight Probe and Bayesian Recommender}\label{sec:probe}

\paragraph{Pre-registered forms (amendments A2-10, A2-11).} For the
pre-registered tests in \S\ref{sec:results}, $\mathrm{ACE}_{\mathrm{norm}}$ is
the entropy of the \emph{direct simplex projection} of the confidence vector
(\texttt{metrics.ace\_d}, no temperature), and the confidence-drop term is the
\emph{unclipped} $1-\Delta C\in[0,2]$; the lifecycle composite renormalises the
two default weights over these two terms and is not clipped, so
$\mathrm{TDI}_{\mathrm{lifecycle}}\in[0,17/12]$. The softmax form above and the
clipped term are reported descriptively beside them. Both amendments were
locked before any v0.6.0 data existed.

\paragraph{Probe signature.} Given a run, the round-$0$ probe signature is
\[\mathbf{x}_0 = (\mathrm{ACE}_{\mathrm{norm}}(0), \overline{\mathbf{C}}(0),
\max_i c_i(0), \mathrm{CSD}(0)) \in \mathbb{R}^4.\]
In practice the probe is
produced by a \emph{shallow} one-round execution of the orchestrator (smaller
tool-call budget, narrower retrieval) so it is materially cheaper than a full
run.

\paragraph{Bayesian recommender.} Let $\Phi$ be a discrete configuration
space of orchestration hyperparameters
$\varphi = (N, R, b)$. We model
$\mathbf{x}_0 \mid \varphi^\star \sim \mathcal{N}(\mu_\varphi,
\Sigma_\varphi)$ with diagonal $\Sigma_\varphi$ and a structured prior
$\pi(\varphi) \propto \exp(-\mathrm{FLOPs}(\varphi) / \kappa)$ that penalises
compute. Given a calibration set
$\{(\mathbf{x}_0^{(k)}, \varphi^{\star(k)})\}_{k=1}^K$, the recommender
picks
\begin{equation}
  \widehat\varphi(\mathbf{x}_0) = \arg\max_{\varphi \in \Phi,\, \mathrm{FLOPs}(\varphi) \leq B}
  \;\; \log \pi(\varphi) + \log \mathcal{N}(\mathbf{x}_0; \mu_\varphi, \Sigma_\varphi).
\end{equation}
With $K\!\sim\!100$ calibration tuples this is cheap, closed-form, and
uncertainty-reportable.

\begin{center}
\fbox{\begin{minipage}{0.92\linewidth}
\textbf{Procedure 1: Bayesian agentic hyperparameter tuning (one task).}\\[2pt]
\textbf{Input:} task $t$, calibrated means $\{\mu_\varphi\}$, variances
$\{\Sigma_\varphi\}$, budget $B$.\\
\begin{enumerate}[leftmargin=2em,topsep=1pt,itemsep=0pt]
  \item $\mathbf{x}_0 \gets \textsc{Signature}(\textsc{ShallowOneRound}(t))$.
  \item For each $\varphi \in \Phi$ with $\mathrm{FLOPs}(\varphi) \le B$:
        compute $\ell_\varphi = \log \pi(\varphi) + \log \mathcal{N}(\mathbf{x}_0; \mu_\varphi, \Sigma_\varphi)$.
  \item Return $\arg\max_\varphi \ell_\varphi$.
\end{enumerate}
\end{minipage}}
\end{center}

% =====================================================================
% Experimental setup for the v0.6.0 sweep. Design facts only; every result
% value is \pending{} (defined in paper.tex).

\section{Experimental Setup}\label{sec:setup_exp}

\subsection{Data and provenance}\label{subsec:datasets}

We use two published K562 Perturb-seq screens in the scPerturb
repackaging~\citep{peidli2024scperturb} (Zenodo record $13350497$):
\textsc{Adamson 2016}~\citep{adamson2016perturbseq}, a CRISPR interference
screen, from its three 10X runs (pilot, 10X005, 10X010) concatenated onto the
intersection of their gene vocabularies; and
\textsc{Norman 2019}~\citep{norman2019perturbseq}, a CRISPR activation screen
with single and double perturbations. All four files are fetched
programmatically; each file's SHA-256 digest is pinned in
\texttt{src/perturb\_eval/data/download.py}, the digests were cross-checked
against the size and MD5 published in the Zenodo record, and a missing pin or
a digest mismatch refuses the fetch. Perturbation labels follow a recorded
label contract. Adamson raw construct labels are parsed structurally: the
plasmid suffix and \texttt{\_only} are stripped; a one-component construct
becomes a gene task (plasmids for the same gene are pooled and the guide
count per gene recorded); multi-gene constructs are excluded; negative-control
constructs are controls; cells with a missing perturbation annotation are
excluded and counted; and labels whose target cannot be corroborated in the
data are excluded with the reason. Norman symbols renamed since publication
are joined to the dataset vocabulary on their Ensembl gene IDs, and labels
that resolve under neither name are excluded with the reason. The held-out
task plan is drawn deterministically: Adamson targets are binned into $3$
quantile bins of mean $|\Delta\log1\mathrm{p}|$ on the target gene and $7$
are drawn per bin from the structurally eligible single-gene constructs
($21$ tasks); Norman contributes $15$ singletons and $5$ doublets ($20$
tasks), stratified by CRC32-derived strata, for $41$ tasks in total, all with
\texttt{seed=2026}. The sampler asserts that every stratum is filled exactly
and the run is refused otherwise. Every excluded label, every digest, the
resolved task plan with its strata and eligible-pool sizes, and the per-label
cell cap (\texttt{max\_cells\_per\_pert}) are written to the run's
provenance record (\S\ref{sec:reproducibility}).

\paragraph{On-target signal.} Each backbone receives the signed
log-fold-change of the perturbation's target gene(s) as its on-target
feature; for a doublet it is the mean over the two target genes. In Adamson
(CRISPR interference) the on-target signal is a decrease; in Norman (CRISPR
activation) it is an increase. Because the feature is signed, the direction
is learned from the data rather than assumed.

\subsection{Features and evaluation}\label{subsec:eval}

\paragraph{Features.} For each held-out task, the top-$2000$ highly variable
genes are selected on that task's \emph{training} cells only, with the target
gene(s) force-included. Models for different tasks therefore use different
feature spaces, and every cross-task median reported here is a median over
models with differing inputs.

\paragraph{Evaluation genes.} The error metric is the mean squared deviation
(MSD) between predicted and observed mean log-fold-change on the $20$ genes
with the largest absolute mean difference between the held-out
perturbation's cells and control cells --- the CPA/GEARS
convention~\citep{lotfollahi2023cpa,roohani2024gears}. These genes are chosen
using the held-out cells.

\subsection{Backbone pool}\label{subsec:backbones}

The Architect picks from three backbones: \textsc{linear} (ridge regression
per gene, NumPy); \textsc{mlp} (a two-layer MLP, PyTorch); and a
$2.1$\,M-parameter gene-token transformer trained from scratch (code
identifier \texttt{scgpt\_small}; not a pretrained foundation model).

\subsection{Trainer sweep and the oracle quantity}\label{subsec:oracle}

For every task the trainer sweep fits all $54$ configurations
$\{\textsc{linear}, \textsc{mlp}, \texttt{scgpt\_small}\} \times
N\in\{3,5\} \times R\in\{1,2,3\} \times 3$ seeds, with the seed passed to the
trainer. The quantity gated in H1/H2 is, per task, the minimum MSD over the
$54$ configurations, then its median over tasks. The minimum is selected on
the same evaluation genes and cells it is scored on, so it is an
\emph{oracle}: an attainable upper bound on what the configuration space can
reach, not an estimate of held-out performance. There is no nested
validation/test split.

\subsection{Agentic lifecycle}\label{subsec:lifecycle}

Five agents --- DataCurator, Literature, Architect, Trainer, Validator ---
produce Pydantic-validated proposals over the following configuration space:

\begin{itemize}
  \item \textbf{DataCurator}: \texttt{hvg\_method}$\in$\{seurat, scanpy\},
  \texttt{hvg\_count}$\in$\{500, 1000, 2000, 5000\},
  \texttt{qc\_mito\_max}, \texttt{split\_strategy}$\in$\{per\_pert\_holdout, unseen\_gene\},
  \texttt{batch\_correction}$\in$\{none, combat, harmony\}.
  \item \textbf{Literature}: BioGPT mechanism retrieval plus PubMed
  E-utilities and STRING-DB neighbour queries; output is a structured pathway
  prior and PPI neighbours for the Architect.
  \item \textbf{Architect}: backbone $\in$\{linear, mlp, scgpt\_small\},
  \texttt{n\_agents}$\in$[2,8], \texttt{n\_rounds}$\in$[1,5],
  \texttt{hvg\_count}, \texttt{learning\_rate}$>0$,
  \texttt{ridge\_lambda}$\ge 0$, \texttt{epochs}$\in$[1,500].
  \item \textbf{Trainer}: \texttt{lr}, \texttt{epochs},
  \texttt{ridge\_lambda}.
  \item \textbf{Validator}: dynamic threshold $\in$[0.02, 0.3] derived
  from the task's probe signature; returns a
  \texttt{StructuredCritique} with \texttt{which\_genes\_failed} and
  \texttt{suggested\_next\_config\_delta}, which the Architect applies in the
  next round.
\end{itemize}

\paragraph{LLM condition.} The agents are served by a rotating pool of
OpenRouter models with failover between pool members. The pool roster is
recorded in provenance, and every lifecycle step records the model that
served it (\texttt{model\_id}) and is labelled \texttt{source} $\in$
\{\texttt{llm}, \texttt{fallback}\}, where \texttt{fallback} marks the
rule-based default the pool returns after an LLM failure. A run with any
fallback step is invalid by construction: it is marked failed and the
analyser refuses to summarise it. Agent-choice statistics are computed over
LLM-sourced steps only. Each task runs the lifecycle with three seeds; the
seed is passed to the trainer and into the LLM cache key. The per-task seed
medians used for H4 and H5 depend on this: in the v0.5.0 artifacts the seed
did not reach the lifecycle run, and the three seeds were identical runs.

\paragraph{What the lifecycle trace records.} Each step records its role,
round, proposal, LLM confidence, \texttt{source} and \texttt{model\_id}. The
five roles do not critique each other, and no round has a winner. The trace
therefore has no critique matrix and no winner sequence, so CSD, WFR and the
four-component TDI of \S\ref{sec:metrics} cannot be computed on this testbed.
H4 and H5 use only the confidence-derived quantities
(\S\ref{sec:tdi-real}).

\subsection{Compute}\label{subsec:compute}

Preflight, trainer sweep and lifecycle sweep run in one process on a single
Modal A100-40GB container with an in-loop budget cap of \$28; the cap is a
limit, not the spend. Every run is appended atomically to
\texttt{trainer\_runs.jsonl} or \texttt{lifecycle\_runs.jsonl}, whose
record~0 is the run's provenance record. Actual spend and GPU-hours:
\pending{actual cost (USD) and GPU-hours from provenance.json}.


% Results for the v0.6.0 sweep. The sweep has not run: this section states the
% pre-registered hypotheses, gates and estimators, and marks every value with
% \pending{} (defined in paper.tex). The v0.5.0 results file
% (v050_results_filled.tex) is no longer input; its values are superseded
% (Appendix~\ref{app:corrections}).

\section{Results}\label{sec:results}

\subsection{Pre-registered hypotheses and gates}\label{subsec:prereg}

The hypotheses, gates and analysis plan below are fixed in
\texttt{paper/PREREGISTRATION.md}; the sweep's provenance record
(\texttt{git\_sha}, \texttt{git\_dirty}) identifies the revision of that file
in force when the sweep ran, and \texttt{summary.json} carries the
pre-registration pin (path, SHA-256, commit) copied from record~0. There are
five gates; each is reported as PASS, FAIL, or UNEVALUATED with the reason, and
the tally is stated out of five. An unevaluable gate (fewer than $3$ tasks, an
undefined statistic, or a diagnostic run) is UNEVALUATED, never FAIL. Every
estimator is implemented in
\texttt{src/perturb\_eval/experiments/preregistered.py} and called from
\texttt{analyse\_v05\_run}; the function named in each subsection below is the
estimator. Bootstrap CIs are percentile intervals over \emph{tasks}
($B=10\,000$ resamples, \texttt{numpy.random.default\_rng(2026)}, 95\%).

\begin{table}[h]
\small
\centering
\begin{tabular}{p{0.05\linewidth}p{0.36\linewidth}p{0.2\linewidth}p{0.27\linewidth}}
\toprule
& Hypothesis & Gate & Outcome \\
\midrule
H1 & Trainer sweep attains low error on Adamson (oracle,
  \S\ref{subsec:oracle}) & median $<0.20$ & \pending{H1} \\
H2 & Trainer sweep attains low error on Norman (oracle) & median $<0.30$
  & \pending{H2} \\
H3 & Architect exercises its backbone choice (LLM-sourced steps)
  & entropy $\ge 0.5$ nats & \pending{H3} \\
H4 & A trace feature ($\mathrm{ACE}_{\mathrm{norm}}$, $1-\Delta C$ or
  $\mathrm{TDI}_{\mathrm{lifecycle}}$) ranks tasks by held-out error
  & Spearman $\rho>0.5$ for $\ge 1$ of the three, within Adamson or within
  Norman & \pending{H4} \\
H5 & Ridge-calibrated $\mathrm{TDI}_{\mathrm{lifecycle}}$ transfers
  Adamson $\to$ Norman & Spearman $\rho>0.4$ & \pending{H5} \\
\bottomrule
\end{tabular}
\caption{Pre-registered hypotheses and gates. Estimators are given in the
subsections below and in \texttt{paper/PREREGISTRATION.md}.}
\label{tab:gates}
\end{table}

\subsection{Attainable error of the trainer sweep (H1, H2)}\label{subsec:h12}

Estimator: per task, the minimum MSD on the top-$20$ DEGs over the $54$
trainer configurations (an oracle, \S\ref{subsec:oracle}); the gate is the
median over tasks in each dataset. Reported with the median: $n$, IQR,
maximum, the fraction of tasks above the threshold, and a bootstrap CI over
tasks; Norman is also split into singletons and doublets
(\texttt{preregistered.h1\_h2\_stats}).

\begin{table}[h]
\small
\centering
\begin{tabular}{lp{0.5\linewidth}}
\toprule
Dataset & Oracle best-of-54 MSD (median; IQR; max; fraction over gate; CI) \\
\midrule
Adamson 2016 ($21$ tasks) & \pending{Adamson oracle MSD summary} \\
Norman 2019 ($15$ singletons, $5$ doublets) & \pending{Norman oracle MSD summary} \\
\bottomrule
\end{tabular}
\caption{Oracle (best of $54$ configurations, selected on the evaluation
split) MSD on the top-$20$ DEGs. Gates: Adamson median $<0.20$, Norman
median $<0.30$.}
\label{tab:headline-msd}
\end{table}

\subsection{Architect configuration choices (H3)}\label{subsec:h3}

Estimator: the Shannon entropy (nats) of the Architect's \texttt{backbone}
picks over all LLM-sourced Architect steps, pooled across tasks, seeds and
rounds ($\ln 3$ ceiling). Reported with it: pick counts per backbone, the
number of LLM-sourced Architect steps, the number of distinct serving models
and a per-\texttt{model\_id} breakdown, and the number of distinct proposals
per role. The gate is \texttt{preregistered.h3}. Backbone-choice entropy: \pending{Architect backbone entropy (nats)
and pick counts}. Per-role proposal diversity and per-model breakdown:
\pending{per-role distinct proposals and per-model\_id counts}.

\subsection{Trace metrics against held-out error (H4)}\label{sec:tdi-real}

\paragraph{Which quantities.} The four-component TDI of
equation~(\ref{eq:tdi}) is \emph{not evaluated}. Two of its components are
structurally undefined for the five-role lifecycle: CSD needs a critique
matrix and WFR needs a per-round winner, and the lifecycle records neither. The
four-component TDI was therefore never computable on this testbed. H4 uses the
components that lifecycle traces do supply. For each run, with $C(r)$ the
confidences of the LLM-sourced steps in round $r$:
$\mathrm{ACE}_{\mathrm{norm}}$ of the final round, the entropy of the direct
simplex projection of $C(\text{last})$ divided by $\ln N$
(\texttt{metrics.ace\_d}, range $[0,1]$); the unclipped $1-\Delta C$ with
$\Delta C=\overline{C}(\text{last})-\overline{C}(\text{first})$
(\texttt{metrics.delta\_mean\_confidence}, for runs with at least two rounds;
range $[0,2]$, so a run whose confidence fell ranks above one whose confidence
stayed flat); and
\begin{equation}
  \mathrm{TDI}_{\mathrm{lifecycle}} = \tfrac{7}{12}\,\mathrm{ACE}_{\mathrm{norm}}
    + \tfrac{5}{12}\,\bigl(1-\Delta C\bigr) \in \bigl[0,\tfrac{17}{12}\bigr],
\end{equation}
the default weights $\alpha=0.35$ and $\gamma=0.25$ renormalised over the
two terms, with no outer clip. $\mathrm{TDI}_{\mathrm{lifecycle}}$ is a score on
$[0,17/12]$, not an index on $[0,1]$. The softmax entropy
(\texttt{metrics.ace\_norm}, $\tau=1$) and the clipped
$1-\min(\max(\Delta C,0),1)$ are reported descriptively beside them. The
softmax was replaced before any data existed (pre-registration amendment~2):
on confidences in $[0,1]$ it spans only about $[0.92,1]$ for five roles, so
rounded LLM confidences would produce ties that decide a rank test. $\mathrm{TDI}_{\mathrm{lifecycle}}$ is a \emph{different quantity}
from the four-component TDI, and its $\rho$ is not comparable with the
v0.4.1 TDI $\rho$. A round with fewer than two LLM-sourced steps leaves the
component undefined for that run. For a one-round run, $\Delta C$ (and so
$1-\Delta C$ and $\mathrm{TDI}_{\mathrm{lifecycle}}$) is undefined rather
than set to $0$. The $\Delta C=0$ convention would score immediate
convergence as maximal difficulty, a built-in bias against H4.
$\mathrm{ACE}_{\mathrm{norm}}$ is unaffected. Undefined values are reported
with the reason and are not imputed
(\texttt{preregistered.per\_run\_components}). A task with no defined value
drops out of that component's $\rho$ and is counted: \pending{per-dataset
exclusion counts (runs undefined, single-round runs, tasks dropped)}.

\paragraph{Estimator.} The unit is the task. Per task, each component and
the lifecycle's final MSD are the medians over the task's three seeds
(\texttt{preregistered.per\_task\_table}). This relies on the seed reaching
the lifecycle run: before that fix, the three seeds were identical. Spearman
$\rho$ is computed \emph{within} Adamson ($21$ tasks) and \emph{within}
Norman ($20$ tasks) separately, with default weights and no in-sample
calibration (\texttt{preregistered.h4}). All six values are reported with
$n$ and a bootstrap CI, and each needs at least $3$ tasks. The gate passes if
any of the six exceeds $0.5$. The $\rho$ pooled over all $41$ tasks is
descriptive only and never gates. The six tests are dependent:
$\mathrm{TDI}_{\mathrm{lifecycle}}$ is a fixed combination of the other two
within each dataset. Let $p$ be the null probability that a single test's
$\hat\rho$ exceeds $0.5$ (the gate thresholds the statistic, not a $p$-value).
Measured by permutation through the pre-registered estimator itself
(\texttt{scripts/local/prereg\_null\_fwer.py}, seed 2026): $p = 1.32\%$ at
$n=20$ and $1.07\%$ at $n=21$; a Monte Carlo of the whole six-test gate gives a
false-positive rate of $2.35\%$ (ACE$_{\mathrm{norm}}$ and $1-\Delta C$
with identical ranks), $4.71\%$ (reversed ranks) and $5.84\%$ (independent),
against $0.0475\%$ for a single pooled test at $n=41$ --- roughly $50\times$ to
$123\times$ higher, the price of testing within datasets. The range is over
these three dependence arms; intermediate negative dependence was not
simulated. No correction was pre-registered; instead all
six $\rho$ are reported with their $n$ whatever the outcome, so a pass carried
by one test is visible in the table.
The dependence is mechanical, not hypothetical: in a committed fixture
(\texttt{tests/test\_preregistered.py}, \texttt{TestH4ReportAllSix}) with
$1-\Delta C$ held constant, $\mathrm{TDI}_{\mathrm{lifecycle}}$ inherits the
ranking of ACE$_{\mathrm{norm}}$, so a single real signal produces \emph{two}
passing tests. This is why the gate is reported as a six-row table rather than
a pass count, and why the false-positive rate is given as a range: its two
endpoints are the two extremes of a dependence between ACE$_{\mathrm{norm}}$
and $1-\Delta C$ that cannot be measured without assuming the structure
under test.

\begin{table}[h]
\small
\centering
\begin{tabular}{p{0.27\linewidth}p{0.2\linewidth}p{0.2\linewidth}p{0.2\linewidth}}
\toprule
Component & Adamson $\rho$ ($n$, CI) & Norman $\rho$ ($n$, CI) & Pooled $\rho$ (descriptive) \\
\midrule
$\mathrm{ACE}_{\mathrm{norm}}$ (final round) & \pending{rho ACE-norm, Adamson}
  & \pending{rho ACE-norm, Norman} & \pending{rho ACE-norm, pooled} \\
$1-\Delta C$ & \pending{rho 1-dC, Adamson} & \pending{rho 1-dC, Norman}
  & \pending{rho 1-dC, pooled} \\
$\mathrm{TDI}_{\mathrm{lifecycle}}$ (default weights)
  & \pending{rho TDI-lifecycle, Adamson} & \pending{rho TDI-lifecycle, Norman}
  & \pending{rho TDI-lifecycle, pooled} \\
\midrule
CSD, WFR & \multicolumn{3}{l}{structurally undefined for the lifecycle (not evaluated)} \\
\bottomrule
\end{tabular}
\caption{Pre-registered gate: $\rho>0.5$ for at least one of the three
computable quantities within Adamson or within Norman (six tests; pooled
$\rho$ descriptive only). The four-component TDI is not evaluated.}
\label{tab:tdi-real}
\end{table}

\subsection{Cross-dataset transfer (H5)}\label{subsec:h5}

Estimator (\texttt{preregistered.h5}): weights for
$\mathrm{ACE}_{\mathrm{norm}}$ and $1-\Delta C$ are fitted by ridge regression
on the Adamson tasks, with ridge parameter $\alpha=1.0$ fixed in advance.
Each feature is standardised with the Adamson mean and s.d., and the target is
the per-task lifecycle MSD. The fitted weights and the Adamson standardisation
are then applied \emph{unchanged} to the Norman tasks, so no Norman value
enters the fit. The gate is Spearman $\rho$ between Norman TDI and Norman
held-out MSD. Transfer $\rho$: \pending{cross-dataset rho with n and CI}.

\subsection{Run record}\label{subsec:runrecord}

Fallback steps (a valid run has none): \pending{fallback-step count}.
Held-out tasks analysed per dataset: \pending{n tasks per dataset from
summary.json}. Exclusions under the label contract:
\pending{excluded labels and cell counts from provenance}.


% =====================================================================
\section{Implementation and Observability Layer}\label{sec:impl}

So far the paper has been algorithmic. This section is deliberately
engineering-heavy because the paper's reviewer question
``\emph{what did you have to contribute to MassGen?}'' deserves a
precise answer. We separate our work into three buckets --- novel
contribution, thin adapter over existing infrastructure, and explicitly-scoped
future work --- and state each with enough detail that a reviewer can locate
the corresponding file in the released code. A fuller design document is
released as \texttt{docs/DESIGN.md} in the project; what follows is a
one-section summary.

\subsection{The trace schema as the decoupling contract}

Our central engineering contribution is a narrow, immutable trace schema
that decouples the orchestrator from the evaluator. A \texttt{RoundTrace}
carries
\[(\mathrm{round\_index}, \mathrm{agent\_names}[N],
\mathbf{C}[N], \mathbf{S}[N,N{-}1], w, \mathrm{consensus\_score},
\mathrm{compute\_tokens});\]
a \texttt{RunTrace} is an ordered tuple of \texttt{RoundTrace}s plus
$(\mathrm{task\_id}, \mathrm{converged}, \mathrm{backbone})$. Both are
\texttt{frozen=True} Python dataclasses, which makes them safe to serialise,
hash, and pass across processes without defensive copying. The schema is
deliberately \emph{numeric-only}: metrics operate on the numbers alone ---
never on the LLM-facing proposal or critique text --- which makes metric
computation costs negligible (\textit{ca.}\ $1$ ms per run) and also
improves privacy posture (LLM content is never consumed by the metric layer).

\begin{table}[ht]
\small
\centering
\resizebox{\linewidth}{!}{%
\begin{tabular}{lll}
\toprule
File & Kind & Status \\
\midrule
\texttt{types.py}, \texttt{metrics.py} & Novel abstraction + metrics & contributed \\
\texttt{probe.py}, \texttt{bayesian.py}, \texttt{calibration.py} & Novel procedures & contributed \\
\texttt{instrumentation.py} & Adapter for CellForge \texttt{ConsensusResult} & contributed \\
\texttt{massgen\_adapter.py} & JSON entrypoints for MassGen skill & contributed \\
\texttt{massgen\_skill\_draft/} (in-repo draft of a MassGen skill) & Ready-to-submit & not yet PR'd \\
LLM-based severity projection (\S\ref{subsec:massgen_bridge}) & Extractor prompt + pure-Python scaffold & not yet calibrated on live traces \\
\bottomrule
\end{tabular}}
\caption{Contribution audit. See \texttt{docs/DESIGN.md} for the full file-by-file breakdown.}
\label{tab:contrib}
\end{table}

\subsection{Adapting CellForge: duck typing, no patching}

The orchestrator (\path{libs/cellforge-agents/}) exposes a result
object, \texttt{ConsensusResult}, whose \path{.rounds[i].proposals} carry
\path{.agent}, \path{.confidence}, \path{.content}, and whose
\path{.rounds[i].critiques} carry \path{.from_agent},
\path{.on_agent}, \path{.severity}. The adapter in
\path{src/perturb_eval/instrumentation.py} is $\approx\!40$ lines of
pure projection: it duck-types those two interfaces, builds a critique
matrix that excludes self-critique, and reconstructs the per-round winner
as $\arg\max_j (c_j - \overline{S}_{\cdot j})$. No orchestrator patching is
required. Any propose-critique-vote system that exposes these two minimal
shapes gets the full evaluator stack for free.

\subsection{MassGen's native observability is rich but differently
shaped}\label{subsec:massgen_bridge}

\textsc{MassGen}~\citep{massgen2025github} already provides a mature
observability layer --- we are not claiming to add telemetry where there
is none. The relevant components are
\texttt{massgen/coordination\_tracker.py}, which records
\texttt{AgentAnswer} (agent\_id, content, timestamp) and \texttt{AgentVote}
(voter\_id, voted\_for, reason, timestamp) dataclasses, and
\texttt{massgen/execution\_trace.py}, which emits structured markdown with
round, vote, and tool-call entries.

Two structural differences prevent a direct wiring:
(i) \textsc{MassGen}'s \texttt{AgentAnswer} does not carry a numeric
confidence; confidence is \emph{implicit} in how many peer voters endorse
the answer.
(ii) \textsc{MassGen}'s \texttt{AgentVote.reason} is a free-text rationale,
not a numeric severity.
\textsc{MassGen}'s own \texttt{CLAUDE.md} explicitly forbids
keyword/regex-based categorisation of such rationales (\emph{``No
keyword/heuristic matching for categorization or similarity''}), so the
projection must use a small LLM-based rater.

We therefore specify \textbf{a two-step extractor}, drafted in the
repository at \texttt{massgen\_skill\_draft/}:
\begin{enumerate}[leftmargin=*,topsep=0pt,itemsep=1pt]
  \item \emph{Vote-share confidence.} For each round, agent $j$'s projected
        confidence is $\hat{c}_j = (\mathrm{votes\_for}_j)/(N{-}1)$. This
        is deterministic and requires no LLM call.
  \item \emph{LLM-rated severity.} For each vote, project the free-text
        \texttt{reason} onto a severity score in $[0,1]$ via a terse rater
        prompt; see \texttt{massgen\_skill\_draft/prompts/severity\_rater.md}.
        The extractor caches by reason string so at most one LLM call per
        distinct reason per session.
\end{enumerate}

The corresponding MassGen skill artifact
(\path{massgen_skill_draft/skill.yaml}) declares two hooks:
\path{on_session_start} runs preflight and \path{on_session_end} runs
evaluate. Both handlers are
$\leq 20$ lines of glue over \path{perturb_eval.massgen_adapter}, which
keeps the skill minimal and audit-friendly.

\paragraph{What we have \emph{not} done.} We have not (a) submitted the
skill to \textsc{MassGen} upstream, (b) calibrated the severity rater
against human annotators on recorded \textsc{MassGen} traces, or (c) validated
the vote-share confidence projection against an alternative projection
(e.g.\ an LLM-rated confidence). Each of these is one incremental
engineering step and we flag them here so the boundary between our paper's
claims and ``plug it in and it works'' is unambiguous.

% =====================================================================
\section{Discussion}\label{sec:discussion}

\paragraph{Why the question is posed as a question.} The title asks whether
agent confidence entropy predicts task difficulty rather than asserting that
it does, so that a null answer is publishable as it stands; a null or failed
outcome on H4/H5 is reported as the answer and is not traded back into a
claim-form title, abstract or contribution list
(\texttt{paper/PREREGISTRATION.md}).

\paragraph{Reading the gates together.} H1/H2 bound what the trainer sweep
can attain (an oracle, \S\ref{subsec:oracle}); they say nothing about whether
the agents find those configurations. H3 asks whether the Architect uses its
backbone choice at all; H4/H5 ask whether the trace metrics carry difficulty
information. On this run H1--H3 pass; H4 passes on the confidence-drop
component ($1-\Delta C$) and on the composite that inherits it, in both
screens, while confidence entropy runs the opposite way to the hypothesis;
and H5 fails: the Adamson-calibrated composite reaches $\rho=\resHFiveRho$ on
Norman against a gate of $\resHFiveThreshold$. Read together: H4 passes on $1-\Delta C$ in both screens and on
$\mathrm{TDI}_{\mathrm{lifecycle}}$, which, given the narrow observed spread of
$\mathrm{ACE}_{\mathrm{norm}}$, ranks nearly identically; the two
$\mathrm{ACE}_{\mathrm{norm}}$ tests fail the pre-registered one-sided gate; and
the Adamson-calibrated composite does not transfer at the pre-registered level.
The title question is answered as posed: on this evidence, confidence entropy
does not predict difficulty in the pre-registered direction. No mechanism is
inferred; the negative point estimates are reported, not interpreted, and a
directional hypothesis about them would be a new pre-registration.

\paragraph{What a single-round probe can and cannot observe.} A single-round
probe by definition cannot observe convergence. It sees only the starting
distribution of agent confidences and critiques, which encodes a snapshot of
prior agreement but not the response-to-critique dynamics that $\Delta C$ and
WFR summarise. This is a structural property of the minimal probe, not an
empirical finding.

\paragraph{Implications for upstream systems.} For
\textsc{CellForge}~\citep{chen2025cellforge} and
\textsc{MassGen}~\citep{massgen2025github}, one recommendation does not depend
on the gate outcomes: \emph{emit the $\{C(r), S(r), w(r)\}$ trace, and the
serving model of every step, as first-class outputs} of every orchestrator
run. Without them neither difficulty metrics nor the attribution of agent
choices to a model is recoverable after the fact (Appendix~\ref{app:corrections}).

% =====================================================================
\section{Limitations and Threats to Validity}\label{sec:limits}

\paragraph{Backbone scope.} \begin{sloppypar}The backbones are small models trained per task
(\S\ref{subsec:backbones}); no pretrained single-cell foundation model is
evaluated. Integrating pretrained \textsc{scGPT}~\citep{cui2024scgpt},
\textsc{Geneformer}~\citep{theodoris2023geneformer}, or
\textsc{scFoundation} is the most consequential follow-up.\end{sloppypar}

\paragraph{Oracle, not held-out, error.} H1/H2 select the best of
\resNDistinctConfigs{} distinct configurations (\resNTrainerRecordsPerTask{}
records per task) on the same evaluation genes and cells they are scored on
(\S\ref{subsec:oracle}). There is no nested validation/test split; the
quantity is an attainable upper bound.

\paragraph{Evaluation genes and feature spaces.} The top-$20$ DEGs are chosen
using the held-out cells (the CPA/GEARS convention), and each task's HVG set
is selected on that task's training cells, so the cross-task medians
aggregate models with different inputs (\S\ref{subsec:eval}).

\paragraph{LLM condition: a single model family.} All five roles are served
by Anthropic Claude models (amendment~4), so nothing here supports a claim of
cross-family generality, and the choice-entropy measurand (H3) may be lower
than a multi-family pool would give. The Validator judges proposals from the
same family; it sits on a different tier (Sonnet 5.5) from the four Haiku 4.5
proposer roles, which is reported, not a remedy --- residual same-family judge
bias is acknowledged. Model identity is recorded and asserted per call.

\paragraph{Sample size and the ACE sign.} The within-screen tests use $n=\resHFourAdaAceN$
and $n=\resHFourNorAceN$ tasks; every $\rho$ carries a bootstrap CI over tasks
(Table~\ref{tab:tdi-real}), and the pre-registered null false-positive rate of
the six-test gate is a few percent. The confidence-entropy point estimates were negative in both screens (the
Adamson CI excludes zero, the Norman CI includes it); that direction was not
pre-registered, generates
no claim, and would need its own pre-registration to test. On Norman the
doublet stratum's oracle median (\resHTwoDoubletMedian) lies above the H2 gate
while the all-task median passes; the stratum is reported, as pre-registered.

\paragraph{Cell-line and dataset scope.} Both datasets are K562, and the
Architect's backbone-choice distribution is K562-conditioned. Genome-scale
validation on \textsc{Replogle 2022}~\citep{replogle2022perturbseq} and a
non-K562 follow-up are explicit next steps.

\paragraph{Configuration space.} $|\Phi|$ remains modest. Widening to
include continuous learning-rate schedules, per-agent retrieval depth,
and explicit batch-correction ablations is straightforward but would
expand the trainer sweep beyond its budget cap. Our Bayesian recommender is
also a MAP estimator; full posterior sampling would expose uncertainty
currently swept into point predictions.

\paragraph{Reward model free-riding.} We deliberately do not train a
process reward model. A PRM over agent traces could outperform hand-designed
metrics, but it would also require an expensive labelling phase. The
question tested here is whether hand-designed difficulty metrics over the
already-produced trace carry the signal; a PRM is a complementary rather
than competing direction.

% =====================================================================
\section{Reproducibility}\label{sec:reproducibility}

The sweep is one Modal command,
\texttt{modal run scripts/modal/app\_v05.py::entrypoint} (with the pinned
\texttt{--prior-spend-usd}), which runs a
CPU preflight, the trainer sweep and the lifecycle sweep in one process.
The preflight fetches and verifies the dataset digests, builds and validates
the task plan, and refuses the whole run --- before any GPU work --- if the
Anthropic key is absent, any roster model is unprobed or any role has fewer than two live models (each model is probed with every role's JSON schema), a pool is smaller than its requested count (the draw is topped
up or trimmed to the per-pool total, A2-7), or a target gene cannot be resolved.

\paragraph{Where each provenance element lives.} ``Provenance-complete'' in
the title refers to the following artifacts; all paths are relative to the
project directory unless stated.
\begin{center}
\small
\begin{tabular}{p{0.38\linewidth}p{0.54\linewidth}}
\toprule
Element & Artifact \\
\midrule
git SHA and dirty flag & \texttt{git\_sha}, \texttt{git\_dirty} in record~0 of
  \texttt{trainer\_runs.jsonl} and \texttt{lifecycle\_runs.jsonl}, and in
  \texttt{provenance.json} \\
resolved entrypoint kwargs & \texttt{entrypoint\_kwargs} in the same records;
  copy in \texttt{configs/runs/<run\_id>.json} \\
resolved library versions & \texttt{lib\_versions} (versions imported at run
  time, not constraint strings) \\
dataset digests & \texttt{datasets} (path, SHA-256, cell and gene counts per
  file) \\
task plan & \texttt{tasks} (resolved lists, strata, eligible-pool sizes) \\
label contract & per-dataset \texttt{label\_contract} and
  \texttt{tasks\_excluded} (label and reason) \\
HVG mode & \texttt{hvg\_selection} (train-only; per-task HVG counts) \\
LLM pool & \texttt{llm\_pool} (roster) and \texttt{llm\_roster\_liveness} (per model: probe date, verdict, per-role result) \\
serving model per step & \texttt{model\_id} and \texttt{source} on every
  lifecycle step in \texttt{lifecycle\_runs.jsonl} \\
run outcome & \texttt{status}, \texttt{cost\_usd\_actual},
  \texttt{gpu\_seconds}, counts in \texttt{provenance.json} \\
hypotheses and analysis plan & \texttt{paper/PREREGISTRATION.md}, at the
  revision named by \texttt{git\_sha} \\
known defects outside the sweep path & repository file
  \path{workstreams/perturb-seq-eval/deferred-findings.md} \\
\bottomrule
\end{tabular}
\end{center}
\begin{sloppypar}
The v0.6.0 run files are written to \path{artifacts/v0.6.0/}: run
\texttt{\resRunId} at \texttt{git\_sha} \resGitSha{} (\texttt{git\_dirty}
\resGitDirty), \texttt{prereg\_version} \resPreregVersion; the run manifest
\path{configs/runs/}\texttt{\resRunId}\path{.json} has SHA-256 prefix \resManifestSha. The
cache namespace and output directory of an earlier aborted attempt were archived
with SHA-256 manifests (prefixes \resArchiveShaCache{} and
\resArchiveShaOutputs), never deleted. The revision at which v0.6.0 landed on main
(\resLandSha) differs from the run's \resGitSha{} in \resPostRunPathCount{} source files, in
code paths the run did not reach (\S\ref{subsec:runrecord}); the landed analyser, re-run on
the committed run files, reproduces every committed value of \texttt{summary.json}.
\end{sloppypar}

\paragraph{Pipeline.}
\begin{sloppypar}
\begin{itemize}[leftmargin=*,topsep=0pt,itemsep=0pt]
  \item \path{src/perturb_eval/data/download.py} --- fetches Adamson's
        three scPerturb subsets and Norman 2019 from Zenodo record
        $13350497$; each file's SHA-256 is pinned and verified, and a
        cached file without a pin or with a mismatching digest is refused.
  \item \path{scripts/modal/app_v05.py} --- preflight, trainer sweep and
        lifecycle sweep on one Modal A100 under the pre-registered spend
        lines; atomic JSONL append per run.
  \item \path{src/perturb_eval/experiments/e_v05_real_traces.py} ---
        consumes both JSONLs and writes \texttt{summary.json}; it refuses
        mismatched task sets, runs with a fallback step, and unfinished or
        failed runs.
\end{itemize}
\end{sloppypar}

\paragraph{Seeds.} The task draw uses CRC32-derived strata and
\texttt{seed=2026}; the per-run seed is passed to the trainer and into the
LLM cache key. The provenance record carries what is needed to re-run the
sweep at the recorded SHA and configuration. Pool-model availability at
re-run time is outside the record's control; the per-step
\texttt{model\_id} states which model served each step.

\paragraph{Software stack.} Pipeline runtime depends on \texttt{numpy},
\texttt{scipy}, \texttt{anndata}, \texttt{h5py}, \texttt{torch} (for
\textsc{mlp} and \texttt{scgpt\_small}), \texttt{pydantic}, \texttt{requests},
and \texttt{modal}; the versions resolved at run time are in
\texttt{lib\_versions}. LaTeX uses \texttt{natbib}, \texttt{booktabs},
\texttt{hyperref}, \texttt{caption}, and \texttt{enumitem}.

\paragraph{Companion design document.} The file \texttt{docs/DESIGN.md} in
the project records the contribution audit and the MassGen bridge
specification.

% =====================================================================
\section{Conclusion}\label{sec:conclusion}

We have proposed ACE, CSD, $\Delta C$-based convergence signals, and WFR as a
four-feature lens on multi-agent group-generation traces, combined into a
Task Difficulty Index, and positioned the TDI and its probe as the
multi-agent analogue of \citet{snell2024scaling}'s compute-optimal
single-chain allocation. We test whether these signals predict task
difficulty under five pre-registered gates on $41$ held-out K562 tasks, with
every run carrying its revision, configuration, dataset digests, task plan,
label contract and per-step serving model. The answer, as pre-registered: \resGateSummary. Within each screen the change in confidence across rounds, and the composite
that inherits it, rank tasks by held-out error at the pre-registered level;
confidence entropy does not, in the pre-registered direction; and the
calibrated composite did not transfer across screens at the pre-registered
level.

\paragraph{Code and data.} All artifacts are released under Apache-2.0 at the
repository that accompanies this submission, including the
\textsc{MassGen}-installable skill that produces per-run metric reports for
any propose-critique-vote orchestrator.

% =====================================================================
\small
\bibliographystyle{plainnat}
\bibliography{references}

\normalsize
\appendix
% Corrections relative to v0.5.0. This is the ONLY place v0.5.0 values
% appear, each marked superseded and mapped to the register row(s) that
% supersede it. Register ids: R*/A* = workstreams/perturb-seq-eval/qa/_adhoc/
% 2026-09-23-publication-rigor-review.md (review + Addendum); DF-* =
% workstreams/perturb-seq-eval/deferred-findings.md; #227 = CTO ruling on HVG
% selection over all cells.

\section{Corrections relative to v0.5.0}\label{app:corrections}

All values reported in the v0.5.0 release of this work are superseded; none
is carried into this version. The v0.5.0 artifact set does not describe a
single experiment, and several values were computed on inputs a code defect
had altered. The defects, by register id (review and addendum:
\texttt{workstreams/perturb-seq-eval/qa/\_adhoc/2026-09-23-publication-rigor-review.md};
DF rows: \texttt{workstreams/perturb-seq-eval/deferred-findings.md}):

\begin{description}[leftmargin=1.2em,itemsep=1pt,font=\normalfont\bfseries]
  \item[A1] The trainer and lifecycle files covered almost disjoint task
    sets (2 shared of 36 each), produced by two processes whose task draws
    diverged (A6: process-salted stratification); statistics joining them
    were computed across that intersection.
  \item[A2] The seed was recorded but not passed to the lifecycle trainer or
    the LLM cache key; the three seeds per task were byte-identical copies.
  \item[A3] Lifecycle steps did not record whether a proposal came from an
    LLM or the rule-based fallback; four of five roles emitted their schema
    default throughout.
  \item[A4] Norman doublet labels were tested with the wrong delimiter, so
    the doublet stratum was empty and each doublet task was trained and
    scored with a randomly chosen gene as its target.
  \item[A7] No dataset digest was pinned; the fetchers were described as
    SHA-gated but trusted any cached file.
  \item[DF-06] Norman labels that did not resolve to a symbol were given a
    random target and entered the training set of every Norman task not
    holding them out.
  \item[DF-07] The Adamson task \texttt{ATF6} pooled four distinct
    constructs (single, two doubles, one triple) into one population.
  \item[DF-10] Adamson cells with a missing annotation were decoded as the
    file's last category; the \texttt{YIPF5} population was about 91\%
    unannotated cells, and both \texttt{YIPF5} and \texttt{ZNF326} entered
    training for other Adamson tasks and were lifecycle held-out tasks.
  \item[\#227] Highly variable genes were selected on all cells, including
    the held-out perturbation's cells.
  \item[R1, R2, R3, R9] Unevaluated gates dropped from the tally; LLM
    condition misstated and unrecorded; documented Norman design contradicted
    by the artifact; provenance record without revision, seed or digests.
\end{description}

\begin{table}[h]
\small
\centering
\begin{tabular}{p{0.44\linewidth}p{0.18\linewidth}p{0.28\linewidth}}
\toprule
Superseded v0.5.0 value or claim & Superseded by & Replacement \\
\midrule
Adamson median best-config MSD $0.147$; Adamson gate PASS
  & A1, DF-07, DF-10, \#227 & H1, \pending{H1} \\
Norman median best-config MSD $0.131$; Norman gate PASS
  & A1, A4, DF-06, \#227 & H2, \pending{H2} \\
Norman $n=15$ described as $15$ singletons $+\,5$ doublets (artifact: $7+8$)
  & R3, A4 & $15+5$ asserted at preflight \\
$1\,944$ trainer cells ($36$ tasks $\times$ $54$)
  & A1, A4 & $41$ tasks $\times$ $54$ \\
$108$ lifecycle runs over $36$ held-out tasks
  & A1, A2 & $41$ tasks $\times$ $3$ seeds, same task set as the trainer \\
Architect backbone entropy $0.36$ nats; \texttt{scgpt\_small} in $91\%$ of
  $138$ picks; HVG-count entropy $0.30$ nats
  & A1, A2, A3, R2, DF-10 & H3, \pending{H3} \\
TDI and component $\rho$ reported as n/a; transfer $\rho$ n/a
  & R1 & H4, H5 \\
``Two of three pre-registered gates pass''
  & R1 & tally out of five \\
``\textsc{Nemotron-30B} throughout''
  & R2 & rotating pool, \texttt{model\_id} per step \\
Spend \$4.04 over 3.06 GPU-hours; \$28 quoted as the envelope
  & A1, R9 & \pending{actual cost and GPU-hours} \\
``SHA-gated'' fetchers & A7 & four digests pinned, fail closed \\
Byte-equivalent re-run from ``the same code revision'' & R9 & re-run at the
  recorded SHA and configuration \\
\bottomrule
\end{tabular}
\caption{One-to-one map from each superseded v0.5.0 value to the register
rows that supersede it.}
\label{tab:corrections}
\end{table}


\end{document}
